% !TEX program = lualatex
% =============================================================================
%  СУПРОВІДНИЙ лист (ДСТУ 4163:2020, п. 5.21): відмітка про наявність
%  додатків обов'язкова; назви додатків — у відмітці, якщо їх немає в тексті.
% =============================================================================
\documentclass[14pt]{dstu4163}

\dstusetup{
  blank      = corner,
  name       = {Товариство з обмеженою відповідальністю «Економічні дослідження та консалтинг»},
  short-name = {(ТОВ «ЕДК»)},
  address    = {вул. Велика Васильківська, 72,\\ м. Київ, 03150},
  phone      = {(044) 123-45-67},
  email      = {office@edk.example},
  code       = {12345678},
  date       = {18.11.2026},
  index      = {01-12/346},
}

% Керівнику структурного підрозділу: юрособа — у називному відмінку,
% посада й ім'я — у давальному (п. 5.15).
\addressee{%
  Науково-дослідний центр «Зразок»

  Начальнику відділу\\
  статистичних даних\\
  Андрію \surname{Гончаренку}}

\begin{document}

\heading{Про надсилання документів}

Надсилаємо документи для розгляду заявки ТОВ «ЕДК» на отримання доступу
до мікроданих статистичної звітності.

% Нумерований перелік (п. 5.21). Одиночний додаток: \attachment{...}
% Брошуровані — без кількості аркушів; > 10 додатків — «згідно з описом».
\begin{attachments}
  \item Заявка на отримання доступу до мікроданих на 2 арк. в 1 прим.
  \item Опис дослідницького проєкту на 5 арк. в 1 прим.
  \item Зобов'язання щодо нерозголошення інформації на 1 арк. в 2 прим.
\end{attachments}

\signature{Директор}{Ігор \surname{Шевчук}}

\executor{Мельник Ольга 044 123 45 68, o.melnyk@edk.example}

\end{document}
